\section{任务型对话}
\label{sec:nlp-dialogue-task}

用户若只问活动何时举行，查通知即可；若要核对自己一行人的报名，答案还取决于身份、查询条件和业务记录，而且这些条件往往分几轮才说清，中途还可能更正。\term{任务型对话}（task-oriented dialogue）系统在多轮交流中帮助用户完成预订、查询一类业务，它面对四个挑战：理解每轮话语的意图与槽位；跨轮累积用户需求，并正确处理更正和撤回；根据当前需求和服务条件选择下一步动作，正确调用真实服务；把服务返回的结果如实转成自然语言，失败时给出可行的下一步。

经典做法是把这四项分给四个模块组成流水线：轮次理解、对话状态追踪、对话策略和回复生成。各模块随后分别神经化，又被合并为单个序列模型端到端生成；大语言模型出现后，状态追踪可以只靠上下文中的示例完成，整个流水线也可以由一个模型配合函数调用实现。本节按流水线的模块顺序展开，每个模块内部从规则到学习方法排列，最后一小节讨论大语言模型如何改变这一分工。

全节使用同一个只读业务。咨询时点为2025年4月17日，用户已登录，身份令牌为U17，目标是核对已有报名而非新增报名。只读接口\code{check\_registration}接受活动编号、认证身份、活动日期和同行人数。假设其快照中有一条U17的记录：活动ML-0418，日期4月18日，人数2，提交于4月16日17:30，状态有效。三轮用户输入为“帮我核对读书会的报名记录”“4月17日，两人”“说错了，是4月18日，人数不变”：第一轮缺日期与人数，第二轮条件下查无匹配，第三轮更正日期后查到记录（图\ref{fig:nlp-dialogue-state}）。

\begin{figure*}[htbp]
  \centering
  % Book-styled flowchart; solid arrows are data, dashed arrows update state.
\begin{tikzpicture}[
  every node/.style={font=\figurefont\small,align=center},
  utterance/.style={book node,text width=38mm,minimum height=16mm,
    fill=FigureInput!10,draw=FigureInput},
  state/.style={book node,text width=42mm,minimum height=16mm,
    fill=FigureModel!10,draw=FigureModel},
  result/.style={book node,text width=48mm,minimum height=16mm,
    fill=FigureOutput!10,draw=FigureOutput},
  flow/.style={book arrow,rounded corners=1.2mm},
  update/.style={flow,dashed,draw=FigureModel},
  labeltext/.style={text=BookMuted,font=\figurefont\small}
]
  \node[labeltext] at (0,14mm) {用户输入};
  \node[labeltext] at (48mm,14mm) {当前需求状态};
  \node[labeltext] at (104mm,14mm) {行为与观察};
  \node[utterance] (u1) at (0,0) {第1轮：核对报名\\日期、人数未给出};
  \node[state] (s1) at (48mm,0) {日期：缺失\\人数：缺失};
  \node[result] (a1) at (104mm,0) {缺必需字段\\询问日期与人数};
  \node[utterance] (u2) at (0,-28mm) {第2轮：新增约束\\4月17日、两人};
  \node[state] (s2) at (48mm,-28mm) {日期：4月17日\\人数：2};
  \node[result] (a2) at (104mm,-28mm) {只读查询：查无匹配\\核对是否指4月18日};
  \node[utterance] (u3) at (0,-56mm) {第3轮：更正日期\\4月18日、人数不变};
  \node[state] (s3) at (48mm,-56mm) {日期：4月18日\\人数：2（保持）};
  \node[result] (a3) at (104mm,-56mm) {只读查询：匹配1条\\核验两人已按时提交};
  \foreach \i in {1,2,3} {
    \draw[update] (u\i.east) -- (s\i.west);
    \draw[flow] (s\i.east) -- (a\i.west);
  }
  \draw[flow] (a1.south) -- (104mm,-14mm)
    -- (0,-14mm) -- (u2.north);
  \draw[flow] (a2.south) -- (104mm,-42mm)
    -- (0,-42mm) -- (u3.north);
  \node[labeltext,fill=white,inner sep=1mm] at (48mm,-14mm) {回复：补全必需条件};
  \node[labeltext,fill=white,inner sep=1mm] at (48mm,-42mm) {回问：让用户确认修改};
  \node[book node,text width=110mm,minimum height=12mm,
    fill=BookAmber!10,draw=BookAmber] (failure) at (53mm,-84mm)
    {独立故障分支：TIMEOUT不是查无记录\\说明失败，按预算重试或回问是否稍后继续};
  \draw[flow,draw=BookAmber] (a2.east) -- (138mm,-28mm)
    -- (138mm,-84mm) -- (failure.east);
  \node[labeltext] at (53mm,-98mm)
    {固定上下文：认证U17、活动ML-0418；不新增或修改报名};
\end{tikzpicture}

  \caption{只读报名核验的三轮对话。用户目标为核对U17在ML-0418的已有报名；实线传递话语、动作和接口结果，虚线表示状态写入或更正。第二轮查无匹配后回问，第三轮只覆盖日期而保持人数。接口故障另走重试或说明失败的分支，不产生报名结论。}
  \label{fig:nlp-dialogue-state}
\end{figure*}
\FloatBarrier

\subsection{轮次理解}
\label{sec:nlp-dialogue-turn}

轮次理解把单轮话语解析为结构化的语义：\term{意图}（intent）是用户这一轮想做的事，如“核对报名”；\term{槽位}（slot）是完成该意图所需的参数，如日期和人数；操作说明本轮对槽位做了什么，如新增、更正或保持。日期解析以给定的咨询时点为参照，省略和指代较多的话语可先按第\ref{sec:nlp-analysis-utterance-rewriting}小节改写成独立句。方法谱系是：意图分类与槽位标注分别训练；JointBERT用共享编码器联合训练两者；模式约束的记录生成进一步输出操作类型和来源轮次。

\subsubsection{意图分类与CRF槽位标注}
\label{sec:nlp-dialogue-intent-crf}

流水线把轮次理解拆为两个已知任务：句级意图用分类器预测，以交叉熵训练；槽位转为BIO标签序列，用线性链CRF标注，以正确路径的负对数似然训练。BIO编码、CRF的路径打分与合法转移掩码见第\ref{sec:nlp-analysis-paradigm-tagging}小节，这里只说明对话中增加的部分。

第二轮按“4月／17日／，／两／人”分词，正确标签为\code{B-DATE I-DATE O B-PARTY I-PARTY}。解码后按分词到字符的映射把相邻标签合并回原文跨度“4月17日”和“两人”，再规范化为2025-04-17和整数2。设意图分类器在“核对报名、取消报名、查询通知”上输出$(0.8,0.1,0.1)$，选择核对报名。意图只决定走哪个业务流程，不代表权限：本节服务在业务层禁用取消动作，意图误判为取消时也不会执行。

这种组合的缺陷在第三轮暴露出来。“说错了，是4月18日”中，普通DATE标签只能标出新日期的位置，无法表示它要替换哪个旧值；“人数不变”没有任何槽位跨度，标注器的输出为空。更正、撤回与保持需要另行定义标签或规则。

\subsubsection{JointBERT联合意图与槽位}
\label{sec:nlp-dialogue-jointbert}

意图和槽位彼此相关：知道意图是“核对报名”，日期更可能是活动日期；反过来，出现日期和人数也提示意图。分别训练两个模型时这种相关无法共享。JointBERT让两项任务共用一个BERT编码器\cite{nlp-dialogue-chen2019}：句首\code{[CLS]}的表示$\bm{h}_0\in\mathbb{R}^{d}$进入意图分类器，每个原词首个子词的表示进入槽位分类器，两项交叉熵相加：
\begin{equation}
\begin{aligned}
 p(c\mid u)&=\operatorname{softmax}(\bm{W}_{c}\bm{h}_0+\bm{b}_{c}),\\
 p(s_i\mid u)&=\operatorname{softmax}(\bm{W}_{s}\bm{h}_{a(i)}+\bm{b}_{s}),\\
 L&=-\log p(c^*\mid u)-\sum_i\log p(s_i^*\mid u).
\end{aligned}
\label{eq:nlp-dialogue-jointbert}
\end{equation}
其中$a(i)$是原词$i$的首子词位置，$\bm{W}_c\in\mathbb{R}^{|\mathcal{C}|\times d}$、$\bm{W}_s\in\mathbb{R}^{|\mathcal{T}|\times d}$分别映射到意图类别和BIO标签。两个头的梯度都回传到共享编码器，这是与分别训练的直接差别。

对第二轮，设正确意图概率为0.8，五个正确标签的概率均为0.9，则$L=-\log0.8-5\log0.9\approx0.750$。基础模型的槽位是逐位置独立分类，可能输出句首\code{I-DATE}这类非法序列；原文也比较了在槽位头上加CRF的版本，以引入标签转移约束。JointBERT与上一小节一样，输出空间仍是单意图加BIO标签，第三轮的“人数不变”需要新增\code{KEEP(party)}一类操作标签才能表达。

\subsubsection{模式约束的语义记录生成}
\label{sec:nlp-dialogue-semantic-record}

要表达更正、撤回和来源，最直接的办法是让模型生成结构化记录。给定当前轮、限定范围的历史和字段模式，模型输出\code{intent}、\code{operation}、\code{slot}、\code{old\_value}、\code{new\_value}和\code{source\_turn}等字段；有标注时以记录词元的负对数似然训练，提示式实现则固定示例和模型版本。解码时用第\ref{sec:nlp-inference-grammar}小节的文法约束屏蔽非法键和枚举值，生成后再用解析器检查类型。

第三轮应生成两条操作：日期\code{REPLACE(2025-04-17, 2025-04-18)}，人数\code{KEEP}；新日期的依据在第3轮，旧日期的依据在第2轮。若上一轮同时讨论了活动日期和报名截止，“改成周五”未指明对象，就应输出待澄清，而不是同时覆盖两个字段。开放值生成省去了为每个取值预设类别，但会出现格式合法、含义错误的记录：原话若是“不是4月18日”，\code{REPLACE(date, 2025-04-18)}键值都合法，含义却正好相反。因此解析后仍要核对来源跨度、否定范围和历史绑定。

\subsection{对话状态追踪}
\label{sec:nlp-dialogue-state}

轮次理解只看一轮，业务调用却需要截至当前的全部有效要求。\term{对话状态追踪}（dialogue state tracking, DST）在每轮之后维护\term{对话状态}（dialogue state）$B_t$，即当前有效的领域—槽位—值集合，而不是历史上出现过的所有词。MultiWOZ提供多领域对话及逐轮状态标注\cite{budzianowski2018multiwoz}，SGD进一步用服务与槽位的自然语言描述定义接口，并区分训练中见过和未见过的服务\scite{rastogi2020sgd}。常用指标是\term{联合目标准确率}（joint goal accuracy），即预测状态与标注状态在所有槽位上完全一致的轮次比例；它要求每轮全部槽位正确，一处错误会在后续各轮持续计错。

槽位值需要区分四种特殊情况：尚未提及为缺失，用户明确说不知道为未知，表示任意值都可接受为不限定，明确拒绝某个值为否定约束。状态追踪有三个挑战：用户更正时要覆盖旧值而非并存；槽位取值开放，无法预先列出所有日期和名称；新服务不断出现，标注数据跟不上。方法谱系为：规则增量更新，可审计但依赖理解结果；TRADE用复制机制生成开放值；SimpleTOD把状态作为整个对话序列的一部分生成；IC-DST用检索到的示例让大语言模型直接输出状态变化，几乎不需要训练数据。

\subsubsection{规则槽位写入与撤回}
\label{sec:nlp-dialogue-state-rule}

规则更新器接收旧状态$B_{t-1}$和本轮解析出的操作集合$\Delta_t$，计算$B_t=U(B_{t-1},\Delta_t)$，并为每个值记录来源轮号。规则来自字段模式与业务含义，不需要训练。表\ref{tab:nlp-dialogue-state-update}列出本节使用的转换；遇到冲突时不按“最后出现者覆盖”处理，而是标记待澄清。

\begin{table}[htbp]
\centering
\begin{tabular}{p{22mm}p{74mm}}
\toprule
操作 & 状态转换\\
\midrule
新增 & 旧槽位缺失时写入值与来源；已有不相容值则待澄清\\
更正 & 明确指向旧槽位时覆盖，旧值移入历史而非有效值\\
保持 & 不改变原值及其来源\\
撤回 & 删除有效值；若只拒绝某值，则写入否定集合\\
未知／不限定 & 分别写入不同标记，不用同一空字符串表示\\
\bottomrule
\end{tabular}
\caption{对话状态的规则更新。缺失是从未提及，撤回是有来源的交流事件，两者在当前有效值表中可能相同，在历史记录中不同。}
\label{tab:nlp-dialogue-state-update}
\end{table}

活动编号和身份由已认证的上下文提供，以下只演算用户可修改的两个槽位：$B_0=B_1=\varnothing$；$B_2=\{(\text{日期},4/17),(\text{人数},2)\}$，两值均来自第2轮；$B_3=\{(\text{日期},4/18),(\text{人数},2)\}$，日期来自第3轮，人数仍来自第2轮。第三轮的变化集只有一个覆盖操作，不能因为本轮没有提到人数就把2清空。增量规则每次改动都可审计，但前面误写的值会一直保留下去，其可靠性完全取决于输入操作是否正确。

\subsubsection{TRADE复制式状态生成}
\label{sec:nlp-dialogue-trade}

规则更新依赖轮次理解的输出，而早期神经状态追踪器又多把每个槽位当作在预定义取值表上的分类，遇到表外的日期或名称就无能为力。TRADE（transferable dialogue state generator）直接读取对话历史，为每个领域—槽位对生成值，生成时可以从历史中复制词元，因此不需要预定义取值表\scite{nlp-dialogue-wu2019}。

模型结构如下。对话历史的词元经嵌入和双向GRU编码为$\bm{H}=[\bm{h}_1,\ldots,\bm{h}_n]\in\mathbb{R}^{n\times d}$。对第$j$个领域—槽位对，以领域嵌入与槽位嵌入之和作为解码器的首个输入，所有槽位共享同一个GRU解码器，各自独立解码。第$k$步解码状态为$\bm{g}_{j,k}$，对历史的注意力为$\bm{\alpha}_{j,k}=\operatorname{softmax}(\bm{H}\bm{g}_{j,k})\in\mathbb{R}^{n}$。输出分布按第\ref{sec:nlp-generation-copy-pointer}小节的复制门混合词表分布与历史复制分布，同一词在历史中多次出现时复制概率相加：
\begin{equation}
 p(y\mid j,k)=
 \gamma_{j,k}\,p_{\mathrm{vocab}}(y\mid j,k)
 + (1-\gamma_{j,k})\sum_{i:x_i=y}\alpha_{j,k,i}.
 \label{eq:nlp-dialogue-trade}
\end{equation}
$\gamma_{j,k}\in(0,1)$由解码状态、当前输入嵌入和注意力摘要$\bm{c}_{j,k}=\bm{H}^\top\bm{\alpha}_{j,k}$经sigmoid计算。另有一个\term{槽门}（slot gate）用首步的$\bm{c}_{j,0}$做三分类：\code{none}表示该槽位无值，\code{dontcare}表示用户不限定，\code{ptr}表示采用解码出的值。训练目标为$L=\alpha L_{\mathrm{gate}}+\beta L_{\mathrm{value}}$，前者是槽门交叉熵，后者是正确值词元的负对数似然，监督来自每轮的完整状态。

第三轮解码日期时，设$\gamma=0.2$，词表对“4月18日”的概率为0.1，历史中对应位置的复制概率为0.7，则混合概率为$0.2\times0.1+0.8\times0.7=0.58$；旧日期“4月17日”若词表概率0.2、复制概率0.1，混合后为0.12，解码选择新日期。人数槽位仍可从第2轮复制“两人”，规范化后保留为2。

TRADE的优点是开放取值和跨领域共享：所有槽位共用解码器，新领域中与旧领域相似的槽位（如不同领域的“日期”）可以零样本迁移。它的弱点同样来自复制：历史中失效的旧值随时可被复制出来，生成完全合法却过期的状态；槽门三类也不足以表达未知、撤回和否定约束，“用户从未提及”与“用户明确撤销”都会映射为\code{none}。此外，每轮都从全部历史重新生成全部槽位，计算量随槽位数和历史长度增长。

\subsubsection{SimpleTOD完整状态序列生成}
\label{sec:nlp-dialogue-simpletod}

TRADE只负责状态，策略和回复仍需其他模块。SimpleTOD把整个流水线重写为一个序列：对话历史$C$、对话状态$B$、数据库结果摘要$D$、系统行为$A$和去词汇化回复$S$依次用边界标记连接，由一个预训练因果语言模型按下一词元负对数似然训练\scite{nlp-dialogue-hosseini2020}。这样做的动机是让各模块共享同一预训练模型的参数和语言知识，把模块间的接口变成序列中的位置。

推断时分两段生成。先由$C$生成$B$，在状态结束标记处停下，解析、校验后实际查询数据库，得到$D$（匹配条目数等摘要）；再以$[C;B;D]$为条件生成$A$和$S$。$D$必须来自真实查询，而不是让模型续写；具体返回行保留在程序侧，供最后回填。原文把数据库段作为可选输入，其端到端最佳设置甚至不读数据库结果；本节保留它，是为了让回复以真实返回为依据。第三轮的教学序列可读作“历史：三轮；状态：日期4/18、人数2；数据库：匹配1条；行为：告知已有有效记录；回复：查到您在[date]的[party]人报名记录”，其中“匹配1条”是调用后的观察，不是人数。

与TRADE相比，SimpleTOD共享更多参数，结构也最简单；代价是状态、动作和回复的错误相互影响，状态生成仍需训练数据，且序列随对话增长而变长。

\subsubsection{上下文学习式状态追踪}
\label{sec:nlp-dialogue-icdst}

以上方法都需要大量逐轮标注的状态，每增加一个服务就要补充标注。大语言模型可以通过第\ref{sec:nlp-inference-icl-shots}小节的上下文学习，只看少量示例完成新任务；困难在于状态是结构化对象，示例又必须与当前轮相似才有效。IC-DST（in-context learning for dialogue state tracking）针对这两点设计了输出格式和示例检索\scite{nlp-dialogue-hu2022icdst}。

第一个设计是把状态写成SQL。服务模式写成\code{CREATE TABLE}语句，每个领域一张表，槽位为列；本轮的状态变化写成一条\code{SELECT}查询，槽位约束写在\code{WHERE}子句中。代码模型在预训练中见过大量SQL，这种格式比自定义的槽值列表更容易被正确生成，也可以用SQL解析器检查合法性。第二个设计是只预测本轮变化：输入是上一轮状态$B_{t-1}$（以文本列出）和最近一对系统与用户话语$(A_t,U_t)$，输出本轮的状态变化$\Delta_t$，完整状态由
\begin{equation}
 B_t=B_{t-1}\oplus\Delta_t
 \label{eq:nlp-dialogue-icdst}
\end{equation}
得到，$\oplus$表示用$\Delta_t$中的槽值覆盖或新增$B_{t-1}$中的对应槽位。这样输入长度不随对话轮数增长，示例也只需包含一轮。代码\ref{code:nlp-dialogue-icdst}是按这一格式写的教学提示。

\begin{codeblock}[language=SQL,label={code:nlp-dialogue-icdst}]{IC-DST格式的状态追踪提示（教学示例）}
CREATE TABLE registration(
  event_id text,   -- 活动编号
  date text,       -- 活动日期
  party_size int   -- 同行人数
);
-- 检索到的示例
-- [context] registration-date: 2025-04-20, registration-party_size: 3
-- [system] 按4月20日、三人未查到记录。
-- Q: [user] 记错了，是4月21日
SELECT * FROM registration WHERE date = '2025-04-21';
-- 当前轮
-- [context] registration-date: 2025-04-17, registration-party_size: 2
-- [system] 按4月17日、两人未查到匹配记录，您要核对4月18日吗？
-- Q: [user] 说错了，是4月18日，人数不变
SELECT * FROM registration WHERE
\end{codeblock}

模型续写出\code{date = '2025-04-18';}，解析为$\Delta_3=\{(\text{日期},4/18)\}$，代入式\eqref{eq:nlp-dialogue-icdst}得到与规则更新相同的$B_3$。

示例选择决定了效果。IC-DST从有标注的样本池中按相似度检索最近的$k$个轮次作为示例，相似度的定义是关键：两个轮次的话语可能用词相近，状态变化却完全不同。因此它以两轮状态变化的重合程度作为相似度标签：分别按槽位集合与槽位—值集合计算两轮变化之间的F1，取两者平均$s=\frac12(F_{\mathrm{slot}}+F_{\mathrm{slot\text{-}value}})$；再从每个样本的近邻中取$s$最高与最低的一部分构成正、负样本对，微调一个句向量检索器，使状态变化相似的轮次在向量空间中接近。例如当前轮的变化只有“日期改为4/18”，候选甲只把日期改为4/21，候选乙同时改日期和人数。槽位F1甲为1，乙为$2\times1\times\frac12/(1+\frac12)\approx0.667$；槽位—值F1两者都为0（值不同）；于是$s$甲为0.5、乙约为0.333，检索器应把甲排在前面。

IC-DST的优点是几乎不需要训练数据（原文以Codex为语言模型，在只用MultiWOZ百分之一到百分之十的标注时超过了此前的少样本方法；换用较小的开源代码模型时则不及这些基线），新增服务只需改写表结构和补充少量示例。缺点是增量更新会累积错误：$\Delta_t$一旦漏掉更正，错误会通过$B_{t-1}$传到所有后续轮次；每轮都要调用大模型，成本和延迟高于小模型；撤回、否定约束同样需要在SQL格式中另作约定。

表\ref{tab:nlp-dialogue-dst-methods}汇总四种状态追踪方法。

\begin{table*}[htbp]
\centering
\small
\begin{tabular}{p{18mm}p{30mm}p{36mm}p{34mm}p{26mm}}
\toprule
方法 & 解决的问题 & 核心机制 & 新增代价 & 监督需求 \\
\midrule
规则更新 & 更正与保持的可审计处理 & 按操作类型增量改写状态 & 依赖轮次理解；旧错持续 & 无需训练 \\
TRADE & 预定义取值表无法覆盖开放值 & 共享解码器、复制门、三类槽门 & 可能复制过期值 & 逐轮完整状态 \\
SimpleTOD & 模块分离、参数不共享 & 状态、查询摘要、行为、回复串成一个序列 & 错误相互影响；序列变长 & 逐轮状态、行为与回复 \\
IC-DST & 新服务缺少标注 & 状态写成SQL、检索相似轮次、只预测变化 & 增量误差累积；大模型调用成本 & 少量示例 \\
\bottomrule
\end{tabular}
\caption{对话状态追踪方法对照。前三种方法的状态由训练好的专用模型或规则产生，IC-DST用冻结的大语言模型加检索示例产生。}
\label{tab:nlp-dialogue-dst-methods}
\end{table*}

\subsection{对话策略与业务调用}
\label{sec:nlp-dialogue-action}

状态完整之后，系统还要决定下一步做什么。\term{对话策略}（dialogue policy）根据当前状态和此前的服务返回，从询问、澄清、查询、告知、结束等动作中选择一个。服务模式规定了字段类型、必需条件和可用动作。规则策略可控，但难以覆盖复杂情形；监督学习从对话语料中模仿动作；强化学习以任务成功为奖励优化长程决策。三种方法共有一个前提：动作被选中不等于已经执行，业务结论只能来自真实返回。

\subsubsection{有限状态与槽位驱动策略}
\label{sec:nlp-dialogue-rule-policy}

规则策略由服务设计者写定：缺必需字段就询问；字段冲突或指代不明就澄清；参数齐备且认证有效才允许查询；取得与目标相符的结果后告知并结束。允许动作掩码把不可用动作的分数设为$-\infty$；由于已过报名截止且服务只读，本例任何轮次都不允许新增报名。

第一轮日期与人数均缺失，动作为“请问您核对哪天、几人的报名记录”。第二轮以\code{(ML-0418, U17, 2025-04-17, 2)}查询，返回\code{OK}且记录为空，策略转为条件核对：“按4月17日、两人未查到匹配记录；通知中的活动日期为4月18日，您要核对这个日期吗？”系统不应自行把日期改成通知日期再报告结果。第三轮以更正后的参数查询，返回一条有效记录，规则另验提交时点$17{:}30<18{:}00$，告知并结束。若返回\code{TIMEOUT}，则没有任何关于记录的结论，可选动作是说明失败、按预算重试或请用户稍后再试。

\subsubsection{监督动作预测与结构化行动生成}
\label{sec:nlp-dialogue-learned-policy}

规则数量增多后，可以从带动作标注的对话中学习策略。分类路线把历史摘要、当前状态和此前服务返回编码为$\bm{v}_t$，以$\operatorname{softmax}(\bm{W}\bm{v}_t)$预测离散动作，最小化交叉熵，参数再由类型化的槽位回填；生成路线以同样的输入直接生成动作名与参数组成的结构记录，以记录词元的负对数似然训练，并约束动作名、字段和枚举值。两条路线都与允许动作掩码配合：第一轮若未加掩码时$\Pr(\text{询问})=0.7$、$\Pr(\text{查询})=0.3$，缺必需字段使查询不可用，掩码后只能询问。生成路线若写出人数\code{"两"}，先由规范化器转为整数2或拒绝，不能依赖服务容忍任意格式。

监督学习只能模仿语料中的动作分布。它不知道哪个动作最终带来任务成功，语料中没出现过的情形（如本例第二轮的空结果）也学不到恰当反应；对话越长，早期动作的后果越难从逐轮标签中体现出来。

\subsubsection{强化学习对话策略}
\label{sec:nlp-dialogue-rl-policy}

对话策略的好坏最终体现在整段对话是否完成任务、用了多少轮，这正是序列决策问题。把对话建模为\term{部分可观测马尔可夫决策过程}（partially observable Markov decision process, POMDP）的思路很早就已提出\cite{nlp-dialogue-levin2000,nlp-dialogue-young2013}：用户的真实目标不可直接观测，系统只能看到对话状态和服务返回。

形式化如下。时刻$t$的状态$s_t=(B_t,D_t)$由对话状态和最近一次服务返回组成；动作$a_t$取自对话行为集合，如\code{request(date)}、\code{confirm(date)}、\code{query}、\code{inform}、\code{bye}；用户回应和服务返回决定下一状态。奖励按任务设计，本节采用教学取值：每轮$-1$，鼓励简短；成功完成目标时$+20$；失败结束时$-10$。策略$\pi_{\bm\theta}(a\mid s)$的目标是最大化期望回报
\begin{equation}
 J(\bm\theta)=\mathbb{E}_{\tau\sim\pi_{\bm\theta}}\Bigl[\sum_{t=1}^{T}\gamma^{t-1}r_t\Bigr],
 \label{eq:nlp-dialogue-rl-objective}
\end{equation}
$\tau$是一整段对话轨迹，$\gamma\in(0,1]$为折扣因子。

这里只用到策略梯度的一个结论：目标对策略参数的梯度可写为$\nabla_{\bm\theta}J=\mathbb{E}\bigl[\sum_t\nabla_{\bm\theta}\log\pi_{\bm\theta}(a_t\mid s_t)\,(G_t-\beta(s_t))\bigr]$，其中$G_t$是从$t$起的回报，基线$\beta(s_t)$不依赖$a_t$，减去它不改变梯度期望、只降低方差\scite{rl1-williams1992-reinforce}。也可以改用深度Q学习估计动作价值，早期端到端任务型对话系统即采用此法\scite{nlp-dialogue-li2017e2e}。策略梯度、价值方法及其成立条件可参考\BookChapterRef{chap:rl-learning}{卷七“未知模型的求解”章}。

用一个两动作的教学例子看梯度方向。第二轮查询为空后，策略在“澄清是否指4月18日”与“告知无记录并结束”之间选择，当前概率各为0.5，取$\gamma=1$。澄清的轨迹从第二轮起依次得到$-1$与$-1+20$，$G_2=18$；结束的轨迹得到$-1-10$，$G_2=-11$。取两者均值$3.5$为基线，优势分别为$+14.5$与$-14.5$。对“澄清”的logit求导，$\nabla\log\pi(\text{澄清})=1-0.5=0.5$，$\nabla\log\pi(\text{结束})=-0.5$，两条轨迹的贡献都是$+7.25$：成功的轨迹抬高澄清，失败的轨迹压低结束，二者方向一致。

强化学习需要大量交互，与真人交互成本太高，因此通常先在\term{用户模拟器}（user simulator）上训练。议程式模拟器为每段对话抽样一个用户目标，包括约束（日期4月18日、人数2）和要询问的信息，再维护一个待说动作的栈：系统询问日期时把“提供日期”压栈，每轮从栈顶弹出若干动作作为用户回应\scite{nlp-dialogue-schatzmann2007}。模拟器还可以加入误差模型，按一定概率替换槽值，模拟理解模块的识别错误，使策略学会确认。

这一方法的问题来自模拟器与奖励。模拟器的行为比真人规整，策略会利用它的规律，换到真人时性能下降；奖励设计不当会被钻空子，例如成功判定只看是否给出了“告知”动作而不核对内容，策略就会学到尽早告知。强化学习策略的优势是能针对任务成功优化长程决策、处理语料中少见的情形，代价是需要可靠的模拟器、可检查的成功判定和大量训练交互。

\subsection{回复生成与交流恢复}
\label{sec:nlp-dialogue-response}

最后一个模块把系统行为和真实返回转为自然语言。这里要区分两种状态：对话状态描述用户当前要什么，业务状态描述服务中实际有什么。用户更正日期只改变前者，不会改动报名记录；回复必须在两者之间建立有依据的对应。模板保证字段来自真实返回，但表达单一；SC-LSTM学习按语义条件生成，表达更自然，需要另外检查遗漏与多余。

\subsubsection{去词汇化模板与槽位回填}
\label{sec:nlp-dialogue-template}

模板把结果类型与语言表达分开。设计者为成功、空结果和故障三类返回各写模板，规定每个占位符取自接口返回还是查询条件。成功模板为“查到[date]的[party]人有效报名，提交于[submitted]”；空结果模板为“按[date]、[party]人未查到匹配记录，请核对条件”；故障模板不包含关于记录是否存在的任何断言。去词汇化指训练与生成时用占位符代替具体槽值，最后再回填。

第二轮的\code{OK, records=[]}选空结果模板，日期和人数取自本次请求；第三轮取得匹配行后，成功模板填入“4月18日”“2”“4月16日17:30”。回填前检查返回行的身份与活动是否与请求一致、所有占位符是否消除；若用户已改变目标而旧的异步结果才返回，应按请求版本丢弃。模板没有训练成本、不会编造字段，但表达固定，覆盖新情形要人工增写，也无法纠正上游的错误状态。

\subsubsection{SC-LSTM语义条件回复}
\label{sec:nlp-dialogue-sclstm}

模板数量随行为与槽位组合增长，表达也显得生硬。语义控制LSTM（semantically controlled LSTM, SC-LSTM）从带对话行为标注的语料中学习“哪种语义用哪种说法”，不需要逐词对齐\scite{nlp-dialogue-wen2015}。输入是一个控制向量$\bm{d}_0$，由行为类型和槽位的独热特征组成，表示待表达的全部语义。它在普通LSTM单元之外增加读取门$\bm{r}_t$，每生成一个词就从$\bm{d}$中“划掉”已表达的部分：
\begin{equation}
\begin{aligned}
 \bm{r}_t&=\sigma(\bm{W}_{wr}\bm{w}_t+
                   \alpha\bm{W}_{hr}\bm{h}_{t-1}),\\
 \bm{d}_t&=\bm{r}_t\odot\bm{d}_{t-1},\\
 \bm{c}_t&=\bm{f}_t\odot\bm{c}_{t-1}
       +\bm{\imath}_t\odot\widetilde{\bm{c}}_t
       +\tanh(\bm{W}_{dc}\bm{d}_t).
\end{aligned}
\label{eq:nlp-dialogue-sclstm}
\end{equation}
$\bm{w}_t$为当前输入词向量，$\bm{\imath}_t$、$\bm{f}_t$为输入门与遗忘门，其余与普通LSTM相同。剩余语义$\bm{d}_t$经$\tanh(\bm{W}_{dc}\bm{d}_t)$进入细胞状态，持续影响后续生成。训练目标为
\begin{equation}
 L=-\sum_{t}\log p(w_t^*\mid w_{<t}^*,\bm{d}_0)+\lVert\bm{d}_T\rVert+\sum_{t=0}^{T-1}\eta\,\xi^{\lVert\bm{d}_{t+1}-\bm{d}_t\rVert},
 \label{eq:nlp-dialogue-sclstm-loss}
\end{equation}
第一项是词元负对数似然；第二项惩罚结束时仍未表达的语义，防止漏说；第三项惩罚相邻步之间$\bm{d}$的剧烈变化，$\xi$取较大的底数，使一次划掉多个语义特征的代价急剧增大。生成时先为每个对话行为过生成20个候选，再按$R=-(F_{\mathrm f}+F_{\mathrm b}+\lambda\,\mathrm{ERR})$重排并保留前5个，其中$F_{\mathrm f},F_{\mathrm b}$是前向与反向生成器给出的代价，$\mathrm{ERR}=(p+q)/N$是槽位错误率，$p,q$为遗漏与多余的槽位数，$N$为应表达的槽位总数。

对第三轮结果，设待表达语义为日期、人数、提交时点，$\bm{d}_0=(1,1,1)$。生成日期附近时读取门为$(0.1,1,1)$，$\bm{d}$变为$(0.1,1,1)$；生成人数时门为$(1,0.1,1)$，$\bm{d}$变为$(0.1,0.1,1)$。若句子到此结束，$\lVert\bm{d}_T\rVert_2\approx1.01$，主要来自未表达的提交时点，训练会推动模型继续生成。读取门是软约束，不保证每个槽位恰好出现一次，因此重排时仍要检查遗漏与重复，并用真实记录回填。

\subsection{基于大语言模型与函数调用的任务型对话}
\label{sec:nlp-dialogue-llm-function}

前面四个模块各需一套标注和模型，新增一个服务就要更新理解标签、状态模式、动作集合和回复模板。经过指令微调的大语言模型已经能够读懂业务说明、追踪多轮内容并写出自然的回复，于是可以把整条流水线交给一个模型：系统提示给出业务规则和可调用函数的说明，模型在对话中决定何时调用哪个函数、填什么参数，执行结果回填到上下文后再生成回复。这里的\term{函数调用}（function calling）指模型输出一个结构化的调用请求，由外部执行器运行并返回结果，模型本身不直接访问服务。

第一步是把槽位和接口写成\term{函数模式}（function schema）。函数模式用JSON Schema描述函数名、用途和参数的类型与约束，必需参数对应原先的必需槽位，枚举与格式约束对应槽位取值范围（代码\ref{code:nlp-dialogue-function-schema}）。用户身份不出现在参数中，由执行器从已认证会话注入，避免模型填写或篡改他人身份。

\begin{codeblock}[label={code:nlp-dialogue-function-schema}]{只读报名核验的函数模式（教学示例）}
{
  "name": "check_registration",
  "description": "只读查询当前登录用户在指定活动的已有报名记录；不能新增或修改报名。",
  "parameters": {
    "type": "object",
    "properties": {
      "event_id":   {"type": "string", "enum": ["ML-0418"]},
      "date":       {"type": "string", "format": "date",
                     "description": "用户所说的活动日期，YYYY-MM-DD"},
      "party_size": {"type": "integer", "minimum": 1,
                     "description": "同行人数"}
    },
    "required": ["event_id", "date", "party_size"]
  }
}
\end{codeblock}

第二步是调用—结果—回复循环（算法\ref{alg:nlp-dialogue-function-loop}，图\ref{fig:nlp-dialogue-function-loop}）。每轮用户话语到来后，模型以系统提示、函数模式和完整对话历史为条件生成：若输出普通文本，就作为回复交给用户；若输出函数调用，执行器先按模式校验参数、检查权限，再调用服务，把返回结果作为一条工具消息追加到历史，模型继续生成，直到给出回复或达到调用次数上限。

\begin{algorithm}[函数调用式任务型对话的单轮处理]
  \label{alg:nlp-dialogue-function-loop}
  \begin{algorithmic}[1]
    \Require 系统提示$P$（业务规则），函数模式集合$\mathcal{F}$，对话历史$H$，用户话语$u$，已认证身份$z$，调用上限$K$
    \Ensure 回复$o$，更新后的历史$H$
    \State $H\gets H\oplus u$
    \For{$k=1,\ldots,K$}
      \State $m\gets\mathrm{LM}(P,\mathcal{F},H)$；$H\gets H\oplus m$
      \If{$m$为普通文本}
        \State \Return $o=m$, $H$
      \EndIf
      \State 解析$m$得函数名$f$与参数$\bm{a}$
      \If{$f\notin\mathcal{F}$或$\bm{a}$不符合$f$的模式或$f$需要未授予的权限}
        \State $\rho\gets$错误说明\Comment{不执行，交给模型改正或询问}
      \Else
        \State $\rho\gets\mathrm{Execute}(f,\bm{a},z)$\Comment{身份由执行器注入}
      \EndIf
      \State $H\gets H\oplus\rho$\Comment{工具消息，含成功、空结果、超时等状态}
    \EndFor
    \State \Return $o=$“暂时无法完成查询”的说明，$H$
  \end{algorithmic}
\end{algorithm}

\begin{figure*}[htbp]
  \centering
  % 第九卷 13.3.5：基于大语言模型与函数调用的任务型对话（跨栏，约160 mm宽）。
% FigureInput：用户；FigureModel：大语言模型；FigureOutput：执行器与服务；虚线框：可选的显式状态记录。
\begin{tikzpicture}[
  x=1mm,y=1mm,
  figure picture,
  user/.style={figure node,draw=FigureInput,fill=FigureInputFill,minimum height=14mm,
    minimum width=16mm,inner xsep=2pt},
  model/.style={figure node,draw=FigureModel,fill=FigureModelFill,minimum height=14mm,
    minimum width=40mm,inner xsep=2pt,line width=.9pt},
  exec/.style={figure node,draw=FigureOutput,fill=FigureOutputFill,minimum height=14mm,
    minimum width=30mm,inner xsep=2pt},
  record/.style={figure node,draw=FigureMuted,fill=white,minimum height=8mm,
    minimum width=56mm,inner xsep=2pt,dash pattern=on 2pt off 1.5pt}
]
  \node[user] (user) at (8,30) {用户};
  \node[model] (llm) at (56,30) {大语言模型\\系统提示＋函数模式＋历史};
  \node[exec] (exe) at (110,30) {执行器\\校验参数与权限\\注入认证身份};
  \node[exec] (svc) at (148,30) {服务\\\texttt{check\_registration}};
  \node[record] (state) at (56,6) {状态记录（可选）：日期、人数及来源轮次};

  \draw[figure flow] ([yshift=3mm]user.east) -- node[figure edge label,above=1pt] {话语} ([yshift=3mm]llm.west);
  \draw[figure flow] ([yshift=-3mm]llm.west) -- node[figure edge label,below=1pt] {回复} ([yshift=-3mm]user.east);
  \draw[figure flow] ([yshift=3mm]llm.east) -- node[figure edge label,above=1pt] {函数调用} ([yshift=3mm]exe.west);
  \draw[figure flow] ([yshift=-3mm]exe.west) -- node[figure edge label,below=1pt] {工具消息} ([yshift=-3mm]llm.east);
  \draw[figure flow] ([yshift=3mm]exe.east) -- node[figure edge label,above=1pt] {调用} ([yshift=3mm]svc.west);
  \draw[figure flow] ([yshift=-3mm]svc.west) -- node[figure edge label,below=1pt] {返回} ([yshift=-3mm]exe.east);
  \draw[figure flow,dash pattern=on 2pt off 1.5pt] (llm.south) -- node[figure note,right=1pt] {写出与核对} (state.north);
\end{tikzpicture}

  \caption{基于大语言模型与函数调用的任务型对话。模型以系统提示、函数模式和对话历史为条件，或直接回复用户，或输出函数调用；执行器校验参数与权限、注入认证身份后调用服务，把返回结果作为工具消息追加到历史，模型再据此回复。图下方的状态记录是可选的显式状态，用于审计与核对。}
  \label{fig:nlp-dialogue-function-loop}
\end{figure*}

用三轮案例走一遍。第一轮，模型发现三个必需参数缺两个，不调用函数，直接询问日期和人数，相当于规则策略的询问动作。第二轮，模型输出调用\code{check\_registration(event\_id="ML-0418", date="2025-04-17", party\_size=2)}，执行器校验通过，返回空记录；模型据此回复“未查到匹配记录，您要核对4月18日吗”。第三轮，模型从历史中读出旧参数，只替换日期后再次调用，返回一条有效记录，模型生成告知。若第二次调用超时，工具消息写明\code{TIMEOUT}，系统提示中的规则要求模型说明失败而不推断记录是否存在。这一过程中，轮次理解、状态追踪和策略都隐含在一次生成里：模型从历史中读取槽值，自己判断参数是否齐备，再决定调用还是询问。

表\ref{tab:nlp-dialogue-pipeline-vs-llm}按模块对照两种实现。

\begin{table*}[htbp]
\centering
\small
\begin{tabular}{p{20mm}p{58mm}p{70mm}}
\toprule
环节 & 模块化流水线 & 大语言模型加函数调用 \\
\midrule
轮次理解 & 意图分类与槽位标注，需逐轮标注 & 隐含在生成中，由模型从话语和历史中读取 \\
状态追踪 & 显式状态$B_t$，规则或专用模型更新 & 状态隐含在上下文中；可另要求模型输出状态记录 \\
对话策略 & 规则、监督或强化学习，在固定动作集上选择 & 由系统提示中的规则和模型判断决定询问、调用或回复 \\
业务调用 & 槽位回填到固定接口 & 模型按函数模式生成参数，执行器校验与执行 \\
回复生成 & 模板或条件生成模型 & 模型依据工具消息直接生成 \\
新服务适配 & 补标注、扩动作集、增写模板，常需重训 & 改写函数模式与系统提示，少量示例 \\
可审计性 & 各模块输出可单独检查 & 只有调用记录天然可查，中间判断需另行记录 \\
典型错误 & 模块误差逐级传递；覆盖不到语料外情形 & 参数幻觉、误读旧值、未调用就声称结果、成本与延迟较高 \\
\bottomrule
\end{tabular}
\caption{模块化流水线与大语言模型加函数调用的对照。后者把四个模块合并到一次生成中，以函数模式取代固定的槽位与动作集合；业务正确性仍由执行器校验和真实返回保证。}
\label{tab:nlp-dialogue-pipeline-vs-llm}
\end{table*}

这种实现减少了模块标注和误差传递，新增服务只需描述新函数，也能应对流水线语料中没有的表达。它的问题集中在状态与执行两处。状态隐含在上下文中，模型可能从历史中读出已被更正的旧日期，或把用户说的人数与通知中的数字混淆；对专门评测的研究发现，大语言模型在显式状态追踪上仍落后于专用模型，但给定正确状态时能较好地把对话引向完成\scite{nlp-dialogue-hudecek2023}。因此实践中常让模型在每轮额外输出一份状态记录（相当于把SimpleTOD的状态段保留下来），或调用专门的状态更新函数，使状态可以检查和回放。执行一侧，模型可能编造参数、未调用就声称查到结果，或在用户与环境都能改变共享状态时失去一致；评测因此需要检查最终的环境状态，而不只是回复文本\scite{nlp-intro-barres2025}。写操作还应要求用户显式确认，并限制模型可调用的函数集合。

模型怎样获得按模式生成调用、读取结果并继续推理的能力，属于工具使用的训练问题，第\ref{chap:nlp-solving-and-execution}章讨论工具调用数据的构造、工具检索与参数生成，以及多步执行中的失败恢复。本节的结论是：函数调用把任务型对话的四个模块压缩为“模型决策、执行器校验、真实返回作依据”的循环，流水线中显式维护的状态和权限检查在新框架中并没有消失，而是转移到了执行器和状态记录中。
